\begin{abstract}
Las Smart Cities requieren conectividad ubicua y resiliente para sistemas de monitoreo en tiempo real, lo cual es desafiante en escenarios de cobertura limitada por redes terrestres. Este artículo explora la integración de Redes No Terrestres (NTN) con Internet de las Cosas (IoT) para habilitar arquitecturas híbridas satélite-terrestre. Se presenta una revisión del estado del arte en NTN-IoT según estándares 3GPP Releases 17-19, se propone una arquitectura de integración multi-conectividad para aplicaciones de monitoreo urbano (tráfico, medio ambiente, infraestructura), y se analizan métricas de rendimiento como latencia, confiabilidad, consumo energético y densidad de conexiones. Los resultados de simulaciones y análisis comparativos demuestran mejoras significativas en cobertura (hasta 100% en áreas remotas) y eficiencia energética respecto a soluciones puramente terrestres. Se discuten desafíos como handover seamless, optimización de recursos y escalabilidad, junto con direcciones futuras hacia 6G. Esta propuesta contribuye a la habilitación de monitoreo inteligente y sostenible en entornos urbanos complejos.
\end{abstract}